% GENERATED FILE. Edit canonical lessons, then run npm run generate:books.
% canonical-source-hash: fnv1a64:c675826842b35900
% canonical-lessons: IT-C267-scusi, IT-C267-destra-sinistra, IT-C267-ripasso-strada

\chapter{Scusi, Per Andare A --- Asking the Way}
\label{ch:it-asking-the-way}
\hlchaptermodality{267}

\begin{chapteropening}
\textbf{By the end of this chapter:} \emph{I can stop a stranger politely, ask the way to a place, and follow an answer with a destra, a sinistra and dritto.}

Complete IT-C267-ripasso-strada: ask the way, follow the answer, check it back and thank.
\end{chapteropening}

\section[scusi, per andare a…?]{scusi, per andare a… — the question that finds the way}

\label{lesson:IT-C267-scusi}



\begin{quote}
Say \textquotedblleft{}to go\textquotedblright{}. (\emph{Andare.}) Say \textquotedblleft{}the station\textquotedblright{}. (\emph{La stazione.})

You are in a town you do not know, and the station is somewhere.
\end{quote}

\subsection*{You'll want to know: scusi}
\textbf{scusi} — \textquotedblleft{}excuse me\textquotedblright{}, said to someone you do not know.

You met \textbf{scusarsi}, \textquotedblleft{}to apologise\textquotedblright{}; \emph{scusi} comes from the plain verb inside it, \emph{scusare}, \textquotedblleft{}to excuse\textquotedblright{}, made a polite request. It is the form for \textbf{Lei}, the polite \textquotedblleft{}you\textquotedblright{} of \emph{come sta}, which a stranger in the street is owed. Its \emph{-i} comes back in the next lesson.

\begin{grammarlens}[title={Grammar Lens: per andare a…}]
\begin{quote}
\textbf{Scusi, per andare alla stazione?} — Excuse me, how do I get to the station?
\end{quote}

Literally, \textquotedblleft{}for going to the station\textquotedblright{}. There is no verb of asking in it at all: \textbf{per} plus the infinitive \emph{andare}, then the place, and your voice rises at the end.

The \emph{a}, \textquotedblleft{}to\textquotedblright{}, joins the article in front of the place, and the two become one word:

\begin{quote}
\textbf{alla} stazione · \textbf{al} museo · \textbf{all'}albergo
\end{quote}

\emph{a} + \emph{la} is \textbf{alla}, \emph{a} + \emph{il} is \textbf{al}, and \emph{a} + \emph{l'}, before a vowel, is \textbf{all'}.

There is a second way to ask, with \emph{dove}, \textquotedblleft{}where\textquotedblright{}: \textbf{Dov'è la stazione?}, \textquotedblleft{}where is the station\textquotedblright{}. Both work. \emph{Per andare a…} asks for the \textbf{way}; \emph{dov'è} asks for the \textbf{place}.
\end{grammarlens}

\begin{cousinweb}
\textbf{scusare} is inherited from Latin \emph{excūsāre} (EtymDB 2.1), built on \emph{ex-}, \textquotedblleft{}out of\textquotedblright{}, and \emph{causa}, \textquotedblleft{}a case at law\textquotedblright{}: to excuse someone was to get them out of the charge (Lewis \& Short, under \emph{excuso}). English \textbf{excuse} is the same word, borrowed through French.
\end{cousinweb}

\subsection*{Your turn}
Say these aloud:

\begin{itemize}
\raggedright
  \item \emph{Scusi, per andare alla stazione?} — and let the voice rise
  \item the same, to the museum — \emph{Scusi, per andare al museo?}
  \item \emph{all'albergo}, then \emph{al museo}, then \emph{alla stazione}
  \item ask for the place instead — \emph{Scusi, dov'è la stazione? È lontano?}
\end{itemize}

\begin{tcolorbox}[breakable,colback=teal!4,colframe=teal!35!black,title={Before you move on}]
Who gets \emph{scusi}? (A \textbf{stranger}, anyone you call \emph{Lei}.) Ask the way to the hotel. (\emph{\textbf{Scusi, per andare all'albergo?}}) What is \emph{a} + \emph{la}? (\emph{\textbf{Alla.}}) What did Latin \emph{excūsāre} mean? (To get someone \textbf{out of a charge}.)
\end{tcolorbox}
\section[a destra, a sinistra]{a destra, a sinistra — following the answer}

\label{lesson:IT-C267-destra-sinistra}



\begin{quote}
Stop a stranger and ask the way to the station. (\emph{Scusi, per andare alla stazione?})

Asking was the easy half. Now the answer comes back, and it comes fast.
\end{quote}

\subsection*{You'll want to know: a destra, a sinistra}
\textbf{a destra} — \textquotedblleft{}on the right, to the right\textquotedblright{}. \textbf{a sinistra} — \textquotedblleft{}on the left, to the left\textquotedblright{}. With \textbf{dritto}, \textquotedblleft{}straight on\textquotedblright{}, which you already have, that is every direction a street can send you.

\begin{quote}
\emph{La stazione è a destra.} — The station is on the right. \emph{La prima a sinistra.} — The first on the left.
\end{quote}

That second one is a whole answer in four words: \emph{la prima}, \textquotedblleft{}the first\textquotedblright{}, is the ordinal you learned in chapter twenty-six, agreeing with the street it leaves unsaid.

\begin{grammarlens}[title={Grammar Lens: vada dritto, giri a destra}]
A stranger answering a \emph{Lei} question answers in the \emph{Lei} form, and the directions are instructions:

\begin{quote}
\textbf{Vada} dritto, \textbf{poi giri} a destra al semaforo. — Go straight on, then turn right at the traffic lights.
\end{quote}

\textbf{giri} is \emph{girare}, \textquotedblleft{}to turn\textquotedblright{}, with the same \emph{-i} as \textbf{scusi}: the polite instruction of an \emph{-are} verb. \textbf{vada} is \emph{andare}'s irregular one, \textquotedblleft{}go\textquotedblright{}. You need to \textbf{understand} these two, not build them; the rest of the answer is words you know: \emph{dritto}, \emph{poi}, \emph{al semaforo}, \emph{all'incrocio}.

A long answer comes in the order you will walk it, and \emph{poi} marks each new piece.
\end{grammarlens}

\begin{cousinweb}
\textbf{destra} goes back to Latin \emph{dextera}, \textquotedblleft{}the right hand\textquotedblright{}, from \emph{dexter}, \textquotedblleft{}right\textquotedblright{} (Lewis \& Short). English \textbf{dexterous} is the same word: skilled, as the right hand is. \textbf{sinistra} is Latin \emph{sinistra}, \textquotedblleft{}the left hand\textquotedblright{}; Lewis \& Short call the origin of \emph{sinister} doubtful. Roman augurs counted the left as the lucky side, but Latin poets already used \emph{sinister} for \textquotedblleft{}unlucky, adverse\textquotedblright{} (Lewis \& Short), and that is the sense English \textbf{sinister} kept. EtymDB 2.1 lists both Italian words but gives neither a direct Latin parent.
\end{cousinweb}

\subsection*{Your turn}
Say these aloud:

\begin{itemize}
\raggedright
  \item \emph{a destra}, \emph{a sinistra}, \emph{dritto}
  \item \emph{La prima a sinistra.} — then \emph{La seconda a destra.}
  \item the stranger's answer — \emph{Vada dritto, poi giri a destra al semaforo.}
  \item repeat the answer back to check it — \emph{Dritto, poi a destra al semaforo?}
\end{itemize}

\begin{tcolorbox}[breakable,colback=teal!4,colframe=teal!35!black,title={Before you move on}]
Which word is \textquotedblleft{}left\textquotedblright{}, and which English word shares its Latin? (\emph{\textbf{Sinistra}} — \textbf{sinister}.) What does \emph{giri} share with \emph{scusi}? (The \textbf{-i} of a polite instruction.) What does \emph{la prima a destra} leave unsaid? (\textbf{The street} — \emph{la prima} agrees with it.)
\end{tcolorbox}
\section[(review)]{Review — ask, listen, check, thank}

\label{lesson:IT-C267-ripasso-strada}



\begin{quote}
Say \textquotedblleft{}on the left\textquotedblright{}, then \textquotedblleft{}on the right\textquotedblright{}. (\emph{A sinistra, a destra.})
\end{quote}

\begin{grammarlens}[title={What you've built}]
Asking the way has four steps, and you now have Italian for every one.

\textbf{Ask}: \emph{Scusi, per andare alla stazione?} \textbf{Listen}: \emph{Vada dritto, poi giri a sinistra all'incrocio.} \textbf{Check}, by saying the answer back as a question: \emph{Dritto, poi a sinistra?} \textbf{Thank}: \emph{Grazie!}

The check is the step travellers skip, and it is the one that saves the walk. A stranger who hears their own directions said back will correct the one piece you got wrong.
\end{grammarlens}

\subsection*{Your turn}
Say these aloud:

\begin{itemize}
\raggedright
  \item ask — \emph{Scusi, per andare al museo?}
  \item hear the answer and say it back — \emph{Vada dritto, poi giri a destra.} … \emph{Dritto, poi a destra?}
  \item thank — \emph{Grazie!}
  \item now be the stranger: send someone to the station — straight on, then the first on the left
\end{itemize}

\begin{tcolorbox}[breakable,colback=teal!4,colframe=teal!35!black,title={Before you move on}]
What are the four steps? (\textbf{Ask}, \textbf{listen}, \textbf{check}, \textbf{thank}.) Which one do travellers skip? (\textbf{The check.}) And if you think a friend is lost, how do you advise them to ask, with a reason? (\emph{Secondo me \textbf{dovresti} chiedere, perché la stazione è lontana.})
\end{tcolorbox}
